\scanpage{12}
2. \uline{Fonctions de Weyl}

% slanted
\margin{$\unsure{K^\infty|L^\infty}, L^{\unsure{\infty(K)|\infty}}_{\unsure{K|k}}$ \unsure{changer|toujours} \unsure{$L$} \unsure{en|sur} $K^\infty_{\unsure{K|k}}$} \uline{Définition 1} \deleted{Soit} \deleted{\unsure{$M$} \unsure{un mesuré} (\unsure{mesure} \unsure{$m$}).} On appelle \uline{fonction de Polya} sur \deleted{\unsure{$M$}} \inserted{$I=]0,\omega[$ $(0<\omega\leq+\infty)$} une fonction $P$ réelle définie sur l'espace $K^\infty(I)$ des fonctions mesurables à support compact, ayant les propriétés suivantes :

$P_1\unsure{)|/}$ $P$ est \uline{symétrique}, i.e. $P(f)=P(g)$ si $f$ et $g$ sont spectralement équivalentes

$P_2)$ $P$ est convexe, i.e. $P(\lambda f+(1-\lambda)g)\leq\lambda P(f)+(1-\lambda)P(g)$ si $0\leq\lambda\leq1$

$P_3)$ \unsure{Pour|Soit} \deleted{tout} compact $K$, \deleted{\unsure{Pour}} et toute suite $(f_i)$ dans $K^\infty_K(K)$ tendant vers une \deleted{\unsure{limite}} limite $f$ en mesure, on a $P(f)\leq\underline{\lim}\,P(f_i)$

En fait, $P_3$ est équivalente à la condition, plus forte en apparence

$P'_3)$ \inserted{Pour \deleted{tout} compact $K$,} $P$ est semi continue inférieurement sur \deleted{tout $K$} pour la topologie faible sur $L^\infty(K)$

Il faut montrer que pour \unsure{tout} réel $a$, \deleted{$P^{-1}([\unsure{a|-\infty},\unsure{-\infty|a}[)$} \inserted{l'ens $A$ des} \deleted{est \unsure{fblt} fermé} $f\in L^\infty(K)$ telles que $P(f)\leq a$ est \unsure{fblt} fermé. \deleted{D'après le th. de \unsure{Banach-Dieudonné}} \inserted{Cet ensemble est convexe en} \deleted{\unsure{on est} \illegible{2}} \deleted{il suffit de prouver que son intersection avec toute} \deleted{partie \unsure{fermée|bornée} de $L^\infty(K)$ est \unsure{fblt} \deleted{\unsure{compacte}} \unsure{fermée}.} vertu de $P_2$, \deleted{\unsure{et pour prouver}} \inserted{et en vertu} de $P_3$ il est fermé pour les suites bornées dans $L^\infty(K)$ qui convergent en mesure. Or un tel ensemble $A$ est déjà \unsure{fblt} fermé, car en vertu du th. de Banach-Dieudonné, on est ramené au cas où $A$ est borné dans $L^\infty(K)$. Il est donc \unsure{fblt} \unsure{relt} \unsure{cpt} dans $L^{\unsure{1|2}}(K)$, \deleted{\unsure{donc}} de plus il est \uline{fermé} dans $L^1(K)$ (pour la top. forte, donc pour la
% archive sheet number: 13

